% Standalone resolution manuscript generated from solution.md.
% Compile with XeLaTeX; author and review metadata are maintained in Markdown.
\documentclass[10pt,a4paper]{article}
\usepackage[margin=24mm,headheight=15pt]{geometry}
\usepackage{fontspec}
\setmainfont{DejaVuSerif.ttf}[BoldFont=DejaVuSerif-Bold.ttf,ItalicFont=DejaVuSerif-Italic.ttf,BoldItalicFont=DejaVuSerif-BoldItalic.ttf]
\setsansfont{DejaVuSans.ttf}[BoldFont=DejaVuSans-Bold.ttf,ItalicFont=DejaVuSans-Oblique.ttf]
\setmonofont{DejaVuSansMono.ttf}
\usepackage{amsmath,amssymb,unicode-math}
\setmathfont{latinmodern-math.otf}
\usepackage{xcolor,microtype,parskip,needspace,fancyhdr}
\usepackage{bookmark}
\usepackage{xurl}
\hypersetup{colorlinks=true,urlcolor=blue,linkcolor=blue,pdftitle={A
real Jordan-curve symbol with a nonreal finite Toeplitz
spectrum},pdfauthor={Matthew J. Colbrook},pdflang={en-GB}}
\urlstyle{same}
\setlength{\emergencystretch}{3em}
\providecommand{\tightlist}{\setlength{\itemsep}{0pt}\setlength{\parskip}{0pt}}
\providecommand{\pandocbounded}[1]{#1}
\setcounter{secnumdepth}{0}
\allowdisplaybreaks[1]
\widowpenalty=10000
\clubpenalty=10000
\pagestyle{fancy}
\fancyhf{}
\fancyhead[L]{\small\sffamily RESOLUTION / SP-06}
\fancyhead[R]{\small\sffamily 11 September 2026}
\fancyfoot[L]{\parbox[b]{0.9\textwidth}{\fontsize{8}{10}\selectfont Independent
Codex-agent verification; no external human peer review or formal
certification.}}
\fancyfoot[R]{\footnotesize\thepage}
\begin{document}
{\raggedright\Large\bfseries A real Jordan-curve symbol with a nonreal
finite Toeplitz spectrum\par}
\vspace{0.4em}
{\normalsize\bfseries Matthew J. Colbrook\par}
{\small Department of Applied Mathematics and Theoretical Physics,
University of Cambridge, Cambridge, United Kingdom\par}
{\small\href{mailto:m.colbrook@damtp.cam.ac.uk}{m.colbrook@damtp.cam.ac.uk}\par}
\vspace{0.6em}
\textbf{Status of this manuscript:} Independently checked negative
resolution (Codex-agent review).\\
\textbf{Target:}
\href{https://github.com/ajt60gaibb/OpenProblemsInNLA/blob/b4123194697bdf6f8f82518c1dd7d6c40a30c2e0/eigenvalues-and-inverse-problems/SP-06/README.md}{SP-06,
original catalog snapshot}.\\
\textbf{Reviewed:} 11 September 2026.

The original proof draft was generated in a ChatGPT conversation. A
separate Codex agent independently checked the complete argument against
the exact target; its
\href{https://github.com/ajt60gaibb/OpenProblemsInNLA/blob/main/references/colbrook-additional-2026-09-11/verification/reviews/SP-06-review.md}{detailed
review} records \textbf{PASS}. The mathematical sections below are
unchanged from the reviewed draft. This is independent agent
verification, not external human peer review or a formal proof
certificate. The
\href{https://github.com/ajt60gaibb/OpenProblemsInNLA/blob/main/references/colbrook-additional-2026-09-11/README.md}{submission
record} preserves the original manuscript and verification history.

\section{The target}\label{the-target}

For a Laurent polynomial \(b(z)=\sum_{k=-r}^s b_kz^k\), with \(r,s\ge1\)
and \(b_{-r}b_s\ne0\), the repository asks whether the existence of a
Jordan curve \(\gamma\subset\mathbb C\setminus\{0\}\) on which \(b\) is
real implies \[
 \operatorname{spec}(T_n(b))\subset\mathbb R
 \quad\hbox{for every }n\ge1,
 \qquad T_n(b)=(b_{i-j})_{i,j=1}^n.
\] This is the finite-section implication retained as conjectural after
the correction to the Shapiro--Štampach article.\footnote{\emph{Open
  Problems in Numerical Linear Algebra}, SP-06, ``A real-valued symbol
  on a Jordan curve and real Toeplitz spectra,'' checked 11 September
  2026;
  \href{https://github.com/ajt60gaibb/OpenProblemsInNLA/blob/main/eigenvalues-and-inverse-problems/SP-06/README.md}{canonical
  entry}.}

\textbf{Theorem 1.} The Laurent polynomial \[
 b(z)=-64z^{-2}+8z^{-1}-128-8z-63z^2-16z^3-z^4
 \tag{1}
\] is real on a star-shaped Jordan curve enclosing zero, but \[
 \operatorname{spec}(T_2(b))=\{-128+8i,-128-8i\}.
 \tag{2}
\] Consequently, the universal finite-spectrum implication in SP-06 is
false.

\section{Construction of a genuine Jordan
curve}\label{construction-of-a-genuine-jordan-curve}

Define \[
 a(z)=8z^{-1}+8z+z^2,
 \qquad F(r,t)=r^2-1+\frac{r^3}{4}\cos t.
 \tag{3}
\] For every real \(t\), \[
 F(1/2,t)=-\frac34+\frac1{32}\cos t<0,
 \qquad F(2,t)=3+2\cos t\ge1.
\] For \(1/2\le r\le2\), \[
 \frac{\partial F}{\partial r}
 =2r+\frac34r^2\cos t
 \ge r\left(2-\frac34r\right)\ge\frac r2>0.
 \tag{4}
\] The intermediate value theorem and strict monotonicity therefore give
a unique root \(r=\rho(t)\in(1/2,2)\) of \(F(r,t)=0\).

The function \(\rho\) is continuous. For example, if \(t_j\to t\),
compactness supplies convergent subsequences of \(\rho(t_j)\); any limit
satisfies \(F(r,t)=0\), so uniqueness forces the limit to equal
\(\rho(t)\). Uniqueness also gives \(2\pi\)-periodicity. Alternatively,
(4) permits the implicit function theorem, yielding a real-analytic
periodic function.

Set \[
 \gamma=\{\rho(t)e^{it}:t\in\mathbb R\}.
 \tag{5}
\] Because \(\rho\) is positive and continuous, this is the image of a
continuous map from the unit circle. If \(\rho(t)e^{it}=\rho(u)e^{iu}\),
equality of arguments gives \(t=u\pmod{2\pi}\). Thus the map is
injective, and \(\gamma\) is a Jordan curve. It avoids zero and bounds
the star-shaped region with radial extent \(\rho(t)\).

\section{Reality of the symbol on the
curve}\label{reality-of-the-symbol-on-the-curve}

For \(z=re^{it}\), direct computation gives \[
 \begin{aligned}
 \operatorname{Im}a(re^{it})
 &=\left(8r-\frac8r+2r^2\cos t\right)\sin t\\
 &=\frac{8\sin t}{r}\,F(r,t).
 \end{aligned}
 \tag{6}
\] Equation (6) vanishes identically when \(r=\rho(t)\). Hence \(a\) is
real everywhere on \(\gamma\).

Now let \[
 b(z)=a(z)-a(z)^2.
 \tag{7}
\] A real polynomial in a real number is real, so \(b\) is also real on
\(\gamma\). Expanding (7) gives exactly (1). Its extreme coefficients
are \(b_{-2}=-64\) and \(b_4=-1\), so the required nonvanishing
condition holds with \(r=2\), \(s=4\).

\section{A nonreal finite section}\label{a-nonreal-finite-section}

Using the repository's indexing convention, \[
 T_2(b)=\begin{pmatrix}b_0&b_{-1}\\b_1&b_0\end{pmatrix}
       =\begin{pmatrix}-128&8\\-8&-128\end{pmatrix}.
 \tag{8}
\] Its characteristic polynomial is \((\lambda+128)^2+64\), proving (2)
and Theorem 1.

The mechanism is that polynomial composition preserves reality on the
curve, whereas finite Toeplitz truncation does not commute with
polynomial functional calculus. In particular, \(T_2(a-a^2)\) is not
required to equal \(T_2(a)-T_2(a)^2\).

\section{Scope and checks}\label{scope-and-checks}

This refutes the assertion about every finite section. It makes no claim
to refute a weaker criterion concerning the limiting Toeplitz spectrum;
the repository explicitly distinguishes that question.

All coefficients, the finite matrix, its spectrum, and the inequalities
establishing the Jordan curve are exact. The verification script also
samples the curve, but the curve's existence, injectivity, and reality
are proved by (3)--(6), not inferred from a numerical plot.

The repository submission records \textbf{Solved}, with a
negative-resolution notice and the independent Codex-agent PASS report
linked above.

\section{Sources}\label{sources}

\begin{enumerate}
\def\labelenumi{\arabic{enumi}.}
\tightlist
\item
  \emph{Open Problems in Numerical Linear Algebra}, SP-06,
  \href{https://github.com/ajt60gaibb/OpenProblemsInNLA/blob/main/eigenvalues-and-inverse-problems/SP-06/README.md}{canonical
  statement}, checked 11 September 2026.
\item
  B. Shapiro and F. Štampach, \emph{Non-Self-Adjoint Toeplitz Matrices
  Whose Principal Submatrices Have Real Spectrum}, \emph{Constructive
  Approximation} 49 (2019), 191--226;
  \href{https://arxiv.org/abs/1702.00741v4}{arXiv:1702.00741, version
  4}, including the appended correction, printed pp.~27--28.
  \href{https://doi.org/10.1007/s00365-022-09614-0}{Correction DOI}.
\item
  D. Giandinoto, \emph{On reality of eigenvalues of banded block
  Toeplitz matrices},
  \href{https://arxiv.org/abs/2411.16266}{arXiv:2411.16266},
  introduction and §2. These references supply the problem's provenance;
  the counterexample above is not attributed to them.
\end{enumerate}
\end{document}
